\documentclass[12pt]{article}
\usepackage{geometry}
\geometry{margin=1in}
\usepackage{setspace}
\usepackage{booktabs}
\usepackage{array}
\usepackage{amsmath}

\setstretch{1.2}

\title{Solutions -- Quiz 3\\ \vspace{8pt} \Large ECON 308 -- International Economic Relations}
\author{Instructor: Muhebullah Karimzada}
\date{October 21,  2025}
\usepackage{comment}

\begin{document}
\maketitle

\section*{Multiple Choice (1 point each)}

\begin{enumerate}
    \item \textbf{(c) Factor endowments} \\
    Differences in trade patterns arise because countries differ in relative factor supplies, not in technology or preferences.

    \item \textbf{(b) Canada exports capital-intensive goods} \\
    The Heckscher--Ohlin theorem predicts a capital-abundant country exports the capital-intensive good.

    \item \textbf{(c) Trade raises the real income of the abundant factor and lowers that of the scarce factor} \\
    According to Stolper--Samuelson, trade benefits the abundant factor and harms the scarce one.

    \item \textbf{(b) Increase output of the labor-intensive good and decrease output of the capital-intensive good} \\
    Rybczynski theorem: increasing a factor increases output of the good intensive in that factor, holding prices fixed.

    \item \textbf{(b) The U.S. imported labor-intensive goods despite being capital-abundant} \\
    The Leontief Paradox contradicted the simple H--O prediction for U.S.\ trade.
\end{enumerate}

\newpage

\section*{Problem (15 points)}

\noindent \textbf{Given:}
\begin{itemize}
    \item Home: $K_H = 300$, $L_H = 300$
    \item Foreign: $K_F = 100$, $L_F = 300$
\end{itemize}

\noindent \textbf{Unit Input Requirements:}

\begin{center}
\begin{tabular}{lcc}
\toprule
\textbf{Good} & \textbf{Capital per unit ($a_K$)} & \textbf{Labor per unit ($a_L$)} \\
\midrule
Computers (C) & 6 & 3 \\
Textiles (T) & 2 & 4 \\
\bottomrule
\end{tabular}
\end{center}

\vspace{0.5cm}

\subsection*{(a) Determine which country is capital- and labor-abundant (3 pts)}

\[
\frac{K_H}{L_H} = \frac{300}{300} = 1, \qquad 
\frac{K_F}{L_F} = \frac{100}{300} = 0.33
\]
Since $\frac{K_H}{L_H} > \frac{K_F}{L_F}$, Home has a relatively larger capital supply. 

\[
\Rightarrow \text{Home is capital-abundant, Foreign is labor-abundant.}
\]

\vspace{0.4cm}

\subsection*{(b) Identify which good is capital-intensive (3 pts)}

Compute capital--labor ratios by sector:
\[
\left( \frac{a_K}{a_L} \right)_C = \frac{6}{3} = 2, \qquad
\left( \frac{a_K}{a_L} \right)_T = \frac{2}{4} = 0.5
\]

Since $2 > 0.5$, Computers use more capital per worker.

\[
\Rightarrow \text{Computers are capital-intensive, Textiles are labor-intensive.}
\]

\vspace{0.4cm}

\subsection*{(c) Predict the trade pattern using Heckscher--Ohlin theorem (4 pts)}

\noindent The Heckscher--Ohlin theorem states that each country exports the good that is intensive in its abundant factor. Therefore:
\[
\text{Home: Exports Computers (C), Imports Textiles (T)}
\]
\[
\text{Foreign: Exports Textiles (T), Imports Computers (C)}
\]

\vspace{0.4cm}

\subsection*{(d) Apply the Stolper--Samuelson theorem to Home (5 pts)}

\textbf{Step 1:} Trade raises the relative price of Home's export, Computers:
\[
\frac{P_C}{P_T} \uparrow
\]

\textbf{Step 2:} Under competitive pricing:
\[
P_C = c_C(w,r), \quad P_T = c_T(w,r)
\]
where $w$ and $r$ are the real wage and return to capital.

\textbf{Step 3:} Stolper--Samuelson states:
\[
\frac{P_C}{P_T} \uparrow \Rightarrow \frac{r}{w} \uparrow
\]
Thus, the real return to capital rises and the real return to labor falls.

\textbf{Step 4:} Economic interpretation:
\begin{itemize}
    \item Capital owners (abundant factor) gain from trade.
    \item Workers (scarce factor) lose in real terms.
\end{itemize}

\[
\Rightarrow \text{In Home: } r \uparrow, \ w \downarrow
\]


\end{document}
